\subsection{Keterbatasan Implementasi}
\label{subsec:keterbatasan-implementasi}

Implementasi pada Tugas Akhir ini memiliki beberapa keterbatasan sebagai berikut:

\begin{enumerate}
	\item Implementasi dilakukan dalam lingkungan komputer lokal, VPS, dan jaringan IOTA Testnet.
	      \begin{itemize}
		      \item Semua \textit{client} dan \textit{Server} PRE dijalankan langsung di komputer lokal.
		      \item VPS menjalankan kluster IPFS dan IOTA \textit{Gas Station}.
		      \item \textit{Smart contract} yang digunakan diunggah ke jaringan IOTA Testnet.
	      \end{itemize}
	\item Kluster IPFS yang digunakan dalam implementasi hanya menggunakan konfigurasi \textit{default} sesuai implementasi DecMed sebelumnya. Konfigurasi tersebut terdiri atas tiga buah \textit{peers}.
	\item Segmentasi data sesuai dengan yang dibahas pada \ref{subsec:rancangan-segmentasi-data-rme} hanya dilakukan hingga tingkat kombinasi \textit{dataset category} dan \textit{function category}, kemudian disimpan sebagai \textit{freetext}.
\end{enumerate}

Keterbatasan tersebut menunjukkan bahwa implementasi masih berada pada tingkat pembuktian konsep. Pengembangan berikutnya perlu memperluas cakupan pengujian keamanan, memperjelas mekanisme pencabutan akses, memperkuat verifikasi bukti delegasi dan \textit{wallet proof}, serta menguji sistem pada lingkungan yang lebih dekat dengan operasional layanan kesehatan.
